% ===== §9 Generative Balance: The Suppression of Power's Freezing =====
\section{Generative Balance: Against the Freezing of Power}
\label{p17:sec:balance}

\noindent
This section is the first of the paper's two summits, and the bearer of its second and central claim. Everything to this point has been diagnosis and method: the forms power's damage takes, the layers through which estrangement is caused, the reconstruction of IR, and the strict limit of what symbolic means can do for a real end. The question now is the constructive one. Given that power is real, ineliminable, and---as the last section's empirical caution established---not self-balancing, what is to be done? The answer is the paper's deepest positive thesis, and it is not the elimination of power but the perpetual generation of a force against its \emph{freezing}. The section builds the thesis in three movements---power's ineliminability, the distinction of static from generative balance, and the redefinition of connection it forces---and then, in its second half, grounds all three in an ontology that shows the thesis to be not an ethical demand laid upon generativity from without, but a law generativity already obeys.

\subsection{Power is the by-product of difference}
\label{p17:sub:balance-difference}

The first movement disposes of two consoling errors at once: the optimistic naturalism that expects power to dissolve of itself, and the radical utopianism that proposes to abolish it altogether. Both founder on a single fact: \emph{power is the by-product of difference, and difference is the source of all generativity.}

Wherever two subjects differ, power arises---not as an intrusion but as an entailment of the difference. Difference of information yields the power of knowledge; difference of resources, the power of the economic; difference of capability, the power of the organisational; difference of narrative facility, the power of the discursive. And difference of attachment yields the subtlest and most intimate power of all: the one who loves more is the one who more fears loss, and the one who more fears loss is the one more readily induced to surrender their freedom. \emph{Love itself produces power.} This last instance is decisive, for it forecloses the optimist's last refuge. One cannot appeal to a sufficiency of love to exempt a bond from power, because love is itself among power's sources. The deeper the attachment, the sharper the asymmetry of fear it can generate.

It follows that power cannot be eliminated without eliminating difference---and difference cannot be eliminated without eliminating the very creativity that makes a relation generative. A relation purged of all difference would be a relation in which nothing new could be generated, for generation, on the holonomic account, is precisely what difference makes possible. The utopian cure is therefore worse than the disease: to abolish power by abolishing difference is to abolish generativity in the same stroke.

\begin{claim}[The ineliminability of power]
Power can neither dissolve of itself nor be abolished. It is the by-product of difference---of information, resources, capability, narrative, and attachment, so that love itself generates power---and difference is the irreducible source of generativity. To eliminate power would require eliminating difference, and to eliminate difference would be to eliminate the creativity it grounds. The optimistic naturalism that awaits power's dissolution and the radical utopianism that would abolish it are both false: the one mistakes power for an accident, the other for a removable evil, when it is the permanent shadow of the very difference without which nothing is generated at all.
\end{claim}

\subsection{Static equality versus generative balance}
\label{p17:sub:balance-static}

If power can be neither awaited-away nor abolished, the only coherent aim is to prevent its \emph{solidification}---and here the section draws its central distinction, between two utterly different conceptions of how a relation might be just.

The first is \emph{static equality}: the design, at some founding moment, of a fair arrangement---an equal division, a balanced contract, a settlement of roles and rights---which is then expected to hold. Its fatal flaw is temporal. Difference does not cease after the founding; power continues to accumulate along its many axes; and the fair arrangement, fixed at the start, is steadily overtaken by the asymmetries that go on forming after it. A static equality, however just at its founding, ossifies---and an ossified arrangement, returning always to the same fixed configuration while the real distribution of power drifts beneath it, is in the geometric language a cycle of \emph{zero holonomy}: motion that generates nothing, a settlement that has become a tomb. Static equality is the political form of thermodynamic equilibrium, and equilibrium, as the cosmological sections will insist (\xref{p17:sec:dissipation}), is death.

The second is \emph{generative balance}: not the achievement of a fair state but the perpetual maintenance of a \emph{mechanism that re-balances}. It does not aim at a distribution and freeze it; it aims at keeping alive the capacity to generate counter-force as fast as power accumulates. Its principle is dynamic---
\[
  \text{Power accumulation} \;\longrightarrow\; \text{counter-force generation}
\]
---power grows, and the generation of counter-power grows with it, so that the system never settles into a fixed distribution but remains in the far-from-equilibrium condition that, for a dissipative structure, is the very condition of life. The ecosystem is the standing image: a forest has dominant species, but no species expands without limit, because predators, competitors, environmental change, and new arrivals continually generate the counter-forces that keep any single dominance from becoming total. The forest is not equal. It is balanced---dynamically, perpetually, never finally.

\begin{claim}[Static equality versus generative balance]
Static equality designs a fair distribution and expects it to hold; but difference continues, power re-accumulates, and the fixed arrangement ossifies into a cycle of zero holonomy---the political form of equilibrium, which is death. Generative balance abandons the fixed distribution for the perpetual maintenance of a re-balancing mechanism: power accumulation provokes counter-force generation, and the system stays alive precisely by never settling. Justice in a generative relation is therefore not a state to be achieved but an activity to be sustained---not the equality of a distribution but the liveliness of the counter-force that keeps any distribution from freezing.
\end{claim}

\subsection{The redefinition of connection}
\label{p17:sub:balance-connection}

This reconception forces a redefinition of \emph{connection} itself, against the common understanding that takes deep connection to mean the dissolution of difference---the merging of two into a frictionless one. That understanding is exactly wrong, and the reason is now visible. If difference were dissolved, power would indeed momentarily vanish---but so would creativity, for difference is the source of both. The frictionless union is not the summit of connection but its death: a static equality of the most intimate kind, generating nothing.

The deeper connection is the opposite: the \emph{preservation} of difference, together with the maintenance of conditions under which the differences may continually correct one another. Connection, on this account, does not abolish the two-ness of the two; it keeps them two, and keeps the channel between them open, so that their difference remains a perpetual source of mutual correction and mutual generation rather than freezing into a relation of command. The task is not to remove difference but to ensure that difference \emph{generates rather than dominates}---that the asymmetry it inevitably produces is continually metabolised into creation instead of crystallising into rule. This is close to the Deleuzean refusal of a final and correct order in favour of a system that retains the capacity to produce new differences, new connections, new lines of flight.\footnote{For the conception of a politics that does not install a permanent correct order but sustains the capacity to generate new difference and new \emph{lignes de fuite}---lines of flight that escape any totalising capture---see \textcite{deleuze1987}; the resonance is with the diagnosis of control as the capture of such lines, not with a vitalist optimism the present account explicitly qualifies in \xref{p17:sub:balance-tragedy}.}

\begin{claim}[Connection as the preservation of difference]
The deepest connection is not the dissolution of difference into a frictionless union---which would abolish, with difference, the creativity difference grounds, and is merely static equality at its most intimate---but the preservation of difference under conditions that let the differences continually correct one another. To connect is not to become one but to remain two with the channel open, so that difference goes on generating rather than freezing into command. The axes one (\xref{p17:sec:asymmetry}: the generation of good value) and two (the suppression of freezing) are thus not opposed but draw on a single fuel---difference---which the one turns to creation and the other guards against crystallisation into rule.
\end{claim}

\subsection{The summit: the isomorphism with the dialectic of desire}
\label{p17:sub:balance-desire}

The thesis as stated so far---generate counter-force perpetually, preserve difference, oppose the freezing of power---reads as a normative recommendation, a thing one \emph{ought} to do. The summit of the section is the claim that it is not, at bottom, a recommendation at all, but the political appearance of an ontological law that generativity already obeys. The argument proceeds by an isomorphism with the dialectic of desire, and the isomorphism must be taken in a strong sense---not as a metaphor or an illuminating parallel, but as two instances of \emph{one underlying structure}.

Recall the structure of desire in the register the series has used. The real cannot be fully taken up into the symbolic; were it fully symbolised---were the object-cause of desire ever completely captured, \emph{das Ding} finally filled---desire would terminate, and with it the subject's life as a desiring being. Desire persists \emph{because} symbolisation necessarily fails, because there is always a remainder of the real that escapes the symbolic net. This necessary failure of capture is not a defect of desire but the very condition of its continuation: complete symbolisation would be the death of desire.\footnote{For \emph{das Ding}, the object-cause of desire, and the constitutive failure of symbolisation that sustains desire rather than defeating it, see \textcite{lacan2006}; and for the prior development of this register across the series, \textcite{paperix}.} Desire, that is, has a double structure: it drives toward symbolisation---it wants to be expressed, recognised, taken up into meaning---and it simultaneously \emph{resists complete symbolisation}, withholding a core that, once captured, would extinguish it. This double movement is what maintains the difference between the real and the symbolic, and it is the maintenance of that difference that keeps desire---and so the subject as a living, generative being---alive.

Now the political structure. Difference, the source of both creativity and power, is to a generative relation what the real is to desire. Were difference ever fully encoded into a fixed power-structure---were the relation's asymmetries completely captured in a frozen and total order---generativity would terminate, for difference would have ceased to be a living source of correction and become a settled structure of rule. Generativity persists \emph{because} difference resists complete encoding, because there is always a remainder---a line of flight---that escapes capture by any power-structure. This necessary failure of total capture is not a defect but the condition of generativity's continuation: complete encoding would be its death. And the counter-force whose perpetual generation the section has urged is nothing other than the political name of difference's resistance to its own complete encoding---the same withholding-from-capture that, in the register of desire, keeps desire alive.

\begin{claim}[The ontological law]
Generative balance is not an external ethical demand laid upon generativity but the political appearance of a law generativity already obeys: \emph{the maintenance of generativity requires that a system hold itself incompletely captured by its own products.} Desire, generating symbols, withholds itself from total symbolisation, on pain of extinction; a generative relation, generating power, must withhold itself from total encoding into a frozen power-structure, on the same pain. The resistance of the real to complete symbolisation and the resistance of difference to complete encoding are not analogues but two instances of one structure. Hence \emph{surechigai} receives its ultimate cause: estrangement is difference, like the drive, refusing to be completely symbolised---two subjects who can never fully encode one another, because complete mutual encoding would be the death of desire and so of the generativity between them. The passing-without-meeting is not the failure of the bond but the sign that it is still alive.
\end{claim}

\subsection{The tension: why counter-force must be \emph{generated} if resistance is automatic}
\label{p17:sub:balance-tension}

The isomorphism raises, at once, an objection sharp enough to undo it if unanswered, and answering it deepens the thesis. If the resistance of difference to encoding is, like the drive's resistance to symbolisation, an \emph{intrinsic and automatic} property---something difference simply does, by its nature---then why has the paper insisted (\xref{p17:sub:practices-empirical}) that counter-force does \emph{not} arise on its own, that power left unopposed accumulates and freezes? The two claims appear to contradict: resistance automatic, yet counter-force needing active generation.

The resolution lies in distinguishing two orders. At the \emph{ontological} level, the resistance is indeed automatic: difference does always throw off a remainder, there is always a line of flight, the materially subordinate always retain the infrapolitical weapons of the weak (\xref{p17:sub:practices-newer}). Total capture, at this level, almost never occurs; some residue always escapes. But the symbolic order develops, over time, \emph{second-order mechanisms for capturing the lines of flight themselves}---disciplinary apparatuses, the metrics of a social-credit logic, the weaponisation of the very interdependence through which connection runs, the commodification that absorbs even resistance and sells it back.\footnote{On the capture of relational and infrastructural interdependence itself, so that the channels of connection become instruments of coercion, see \textcite{farrell2019}; the broader point---that an order evolves means to capture the escapes that would otherwise limit it---generalises the diagnosis of control in \textcite{deleuze1987}.} These second-order mechanisms do not abolish the ontological resistance---they cannot---but they suppress its \emph{effect}, capturing each line of flight as it forms.

\begin{claim}[Guarding, not inventing, the lines of flight]
The resistance of difference to total encoding is ontologically automatic, but the symbolic order evolves second-order mechanisms that capture the lines of flight as fast as they form. Hence ``generating counter-force'' does not mean conjuring resistance from nothing---the resistance is always already there---but \emph{guarding the always-present resistance against the second-order mechanisms that would capture it.} It is the protection of the lines of flight, not their invention. This dissolves the tension and unifies the two practical axes the praxis will develop: \emph{non-action} (\emph{wu-wei}) is right toward the ontological resistance, which arises of itself and need not be forced; \emph{action} is required against the second-order capture, which does not relent of itself and must be perpetually opposed.
\end{claim}

\subsection{The discipline: isomorphism, not reduction; and the tragic remainder}
\label{p17:sub:balance-tragedy}

Two disciplines guard the summit against misreading, and the section closes by stating them, for without them the isomorphism would overreach in exactly the ways its critics would expect.

The first is that the isomorphism is \emph{formal, not reductive}. The claim is that the resistance-to-capture in desire and the resistance-to-capture in generative relation share an underlying structure---a double movement toward and away from complete capture---not that political struggle ``is'' desire, nor that the relational field is a libidinal economy in disguise. The series has held this discipline at each of its formal couplings---with predictive coding, with the geometry of phase, with the dynamics of the quantum---and it holds it here: the dialectic of desire furnishes a fully theorised \emph{model} of the structure, and the political thesis is shown to instantiate that structure, without either being collapsed into the other.

The second is that the tragic dimension of the desire-structure must be \emph{kept}, not sweetened. Desire's persistence is also its perpetual lack: to desire is to be marked by a want that no object fills, and the continuation of desire is the continuation of that want. The political isomorph inherits the cost. Generative balance is not a happy equilibrium of perpetual liveliness; it is the refusal of a rest that would be death, purchased at the price of a relation that can never arrive at the stillness of a final fairness, a complete concordance, a last reconciliation. This is where the section rejoins the paper's first claim, and gives it its deepest ground. \emph{Surechigai} is not an accident to be one day overcome but the standing price of remaining alive: the cost difference pays for refusing to be wholly encoded, the cost desire pays for refusing to be wholly symbolised, the cost a generative relation pays for refusing the frozen peace of equilibrium. The passing-without-meeting and the perpetual want are one structure, seen in two registers. The bond that has stopped passing has not been perfected; it has died.

\medskip
\noindent
A final word marks the summit's reach and its restraint. The law stated here---that the maintenance of generativity requires a system to hold itself incompletely captured by its own products---is stated for the intimate bond, and the paper claims it only there. But it plainly does not stop there. Raised to the level of a political body, it yields a criterion for evaluating any such body: whether, in generating its value, it also generates the counter-force that keeps its own power from freezing---whether its growth returns to its creators or is captured, whether its lines of flight are guarded or absorbed. That criterion, and its relation to the theories of non-domination, of agonistic plurality, of the right to justification, and of generative justice, is the seed of a separate work \parencite{eglash2016}, and the paper marks it as such and does not pursue it here. It belongs to the next question---how love becomes free within history---to which, after the positive summit of asymmetry, the paper turns.
